\chapter{Case Analysis: Public Transport in Fergana City}
\label{ch:case}
\label{ch:analysis}

This chapter analyses the case city of the thesis. It describes the
city of Fergana, its public transport system, the main problems faced
by passengers, and the requirements for a digital information service.
The chapter creates the practical context for the system that is
designed and built in Chapter~\ref{ch:empirical}.

\section{City Profile and Transport Context}
\label{sec:city-profile}

Fergana is a regional capital in the east of Uzbekistan with a
population of approximately 300{,}000 people \cite{stat_uz_fergana}.
The city is the administrative and economic centre of the Fergana
Valley, one of the most densely populated regions in Central Asia. Daily
mobility in the city depends mostly on public buses and shared taxis.
Private cars are common but not the main mode of transport for
students, pensioners, and low-income groups.

Public bus transport in Fergana is operated by several local carriers.
There is no single official operator that publishes a complete schedule
or a digital map of routes. Information about routes is shared mostly
through signs on buses, word of mouth, and unofficial posts in social
networks. General map services such as Yandex Maps cover only a part of
the local routes and do not provide reliable real-time data.

\section{Identified Problems}
\label{sec:problems}

During informal observation and conversations with local residents, the
author identified several recurring problems with the current state of
public transport information in Fergana. These problems are
summarised in Table~\ref{tab:problems}.

\begin{table}[H]
\centering
\caption{Main problems of public transport information in Fergana}
\label{tab:problems}
\begin{tabular}{p{0.30\linewidth} p{0.62\linewidth}}
\hline
\textbf{Problem} & \textbf{Description} \\
\hline
No official application & The city has no official mobile or web
application for public transport. \\
Incomplete map data & Yandex Maps covers only some routes; many local
buses are missing. \\
No published schedules & Bus arrival times are not published in any
digital form. \\
Language barrier for visitors & Stop names and route information are
mostly in Russian and Uzbek, which is hard for tourists. \\
Difficulty for new residents & Students and newcomers do not know which
bus goes where. \\
Low digital skills in some groups & Older users find complex
applications difficult. \\
\hline
\end{tabular}
\end{table}

These problems show that the gap is not in the transport itself but in
the \emph{information layer} around it. Buses exist and operate every
day, but the knowledge of routes and stops is not stored in a single,
accessible digital place.

\section{Why Telegram}
\label{sec:why-telegram}

Telegram is one of the most used messengers in Uzbekistan. According to
public reports, Telegram is among the top platforms by daily active
users in the country, together with Instagram and WhatsApp
\cite{datareportal_uz_2024}. Many local businesses, news channels, and
even government services already use Telegram bots and channels to
reach citizens.

For a public transport information system, Telegram has several
practical advantages:

\begin{itemize}
  \item \textbf{Low entry barrier.} Users do not need to install a new
  application. They already have Telegram on their phones.
  \item \textbf{Low data usage.} A text-based bot uses very little
  mobile data, which is important for users with limited internet
  plans.
  \item \textbf{Works on weak devices.} Telegram works well on older
  Android phones, which are common in the region.
  \item \textbf{Free Bot API.} The Telegram Bot API is free and well
  documented \cite{telegram_botapi}, which removes the cost barrier
  for a student project.
  \item \textbf{Familiar interface.} Users already know how to send
  messages, which reduces the need for training.
\end{itemize}

A native mobile application (Android/iOS) could also solve the same
problem, but it would require app store accounts, separate code bases
for two platforms, and a higher technical skill from users. A web
application would require users to remember a URL and open a browser.
A Telegram bot avoids these issues.

\section{Stakeholders and User Groups}
\label{sec:stakeholders}

The system has several types of users. They are described in
Table~\ref{tab:users}.

\begin{table}[H]
\centering
\caption{Main user groups of the proposed system}
\label{tab:users}
\begin{tabular}{p{0.25\linewidth} p{0.30\linewidth} p{0.37\linewidth}}
\hline
\textbf{User group} & \textbf{Main need} & \textbf{Typical query} \\
\hline
Local residents & Daily commute information & ``Which bus goes from
home to the market?'' \\
Students & Routes to university and dormitories & ``How to get from
Istiqlol to the city centre?'' \\
Newcomers and visitors & Basic orientation in the city & ``What buses
stop near the airport?'' \\
Older users & Simple, text-based help & ``Show me the stops of bus 1.''
\\
\hline
\end{tabular}
\end{table}

All groups share one common need: a \emph{simple way to ask a question
in text form and get a clear answer}. This need is the main driver of
the system design in Chapter~\ref{ch:empirical}.

\section{Functional Requirements}
\label{sec:requirements}

Based on the problems and user groups above, the following functional
requirements were defined for the system. They are listed in
Table~\ref{tab:reqs}.

\begin{table}[H]
\centering
\caption{Functional requirements of the Fergana Transport Bot}
\label{tab:reqs}
\label{tab:fr}  
\begin{tabular}{p{0.10\linewidth} p{0.32\linewidth} p{0.50\linewidth}}
\hline
\textbf{ID} & \textbf{Requirement} & \textbf{Description} \\
\hline
FR-1 & Show route stops & The bot must show the list of stops for a
selected bus route. \\
FR-2 & Search route between stops & The bot must find a route between
two stop names entered by the user. \\
FR-3 & Provide taxi information & The bot must show contacts of local
taxi services as an alternative. \\
FR-4 & Collect user suggestions & The bot must accept new transport
information from users for future updates. \\
FR-5 & Help and command list & The bot must show a list of commands
and short usage examples. \\
FR-6 & 24/7 availability & The bot must be hosted on a cloud platform
and respond at any time. \\
\hline
\end{tabular}
\end{table}

\section{Non-Functional Requirements}
\label{sec:nfr}

In addition to functional requirements, the system must meet several
non-functional requirements:

\begin{itemize}
  \item \textbf{Usability.} Commands must be short and easy to remember.
  Examples must be shown in the help message.
  \item \textbf{Performance.} The bot must answer in less than two
  seconds under normal load.
  \item \textbf{Reliability.} The bot must restart automatically after
  a server reboot.
  \item \textbf{Low cost.} The system must run on a free or low-cost
  hosting plan to stay sustainable for a student project.
  \item \textbf{Maintainability.} Route and stop data must be stored in
  a database that can be updated without changing the code.
\end{itemize}

\section{Chapter Summary}
\label{sec:ch2-summary}

This chapter described the case city of Fergana and the current state
of its public transport information. The city has no official digital
information service. General map applications cover only part of the
routes. Several user groups -- residents, students, visitors, and
older people -- need a simple way to get transport information in text
form.

Telegram was chosen as the platform because it is widely used in
Uzbekistan, has a free Bot API, low data usage, and a familiar
interface. Based on the identified problems and user groups, six
functional and five non-functional requirements were defined. These
requirements form the basis of the system design and implementation,
which are presented in Chapter~\ref{ch:empirical}.